\IfFileExists{stacks-project.cls}{%
\documentclass{stacks-project}
}{%
\documentclass{amsart}
}
\usepackage{amssymb}

% For dealing with references we use the comment environment
\usepackage{verbatim}
\usepackage{fancyvrb}
\usepackage{pmboxdraw}
\newenvironment{reference}{\comment}{\endcomment}
%\newenvironment{reference}{}{}
\newenvironment{slogan}{\comment}{\endcomment}
\newenvironment{history}{\comment}{\endcomment}

% For source figures retained as assets
\usepackage{graphicx}

% For commutative diagrams we use Xy-pic
\usepackage[all]{xy}

% We use 2cell for 2-commutative diagrams.
\xyoption{2cell}
\UseAllTwocells

% We use multicol for the list of chapters between chapters
\usepackage{multicol}

% This is generally recommended for better output
\usepackage{lmodern}
\usepackage[T1]{fontenc}
\usepackage{textalpha}

% For cross-file-references
% We prefer the snapshot tmp/xref/#2.aux,
% which latexmkrc replaces atomically, so
% that documents compiled simultaneously
% never read a half-written .aux file.
% The last argument keeps links to #2.pdf.
\usepackage{xr-hyper}
\newcommand{\maybeexternaldocument}[2]{%
  \IfFileExists{#2.aux}{%
    \externaldocument[#1]{#2}%
  }{%
    \PackageWarning{stack}{Missing #2.aux; external references with prefix #1 unavailable}%
  }%
}
\let\externaldocumentorig\externaldocument
\renewcommand{\externaldocument}[2][]{%
  \IfFileExists{tmp/xref/#2.aux}{%
    \externaldocumentorig[#1]{tmp/xref/#2}[#2.pdf]%
  }{%
  \IfFileExists{#2.aux}{%
    \externaldocumentorig[#1]{#2}%
  }{%
    \PackageWarning{stack}{Missing #2.aux; external references with prefix #1 unavailable}%
  }%
  }%
}

% Package for hypertext links:
\usepackage{hyperref}

% For any local file, say "hello.tex" you want to link to please
% use \externaldocument[hello-]{hello}
\externaldocument[introduction-]{introduction}
\externaldocument[conventions-]{conventions}
\externaldocument[sets-]{sets}
\externaldocument[categories-]{categories}
\externaldocument[topology-]{topology}
\externaldocument[sheaves-]{sheaves}
\externaldocument[sites-]{sites}
\externaldocument[stacks-]{stacks}
\externaldocument[fields-]{fields}
\externaldocument[algebra-]{algebra}
\externaldocument[brauer-]{brauer}
\externaldocument[homology-]{homology}
\externaldocument[derived-]{derived}
\externaldocument[simplicial-]{simplicial}
\externaldocument[more-algebra-]{more-algebra}
\externaldocument[smoothing-]{smoothing}
\externaldocument[modules-]{modules}
\externaldocument[sites-modules-]{sites-modules}
\externaldocument[injectives-]{injectives}
\externaldocument[cohomology-]{cohomology}
\externaldocument[sites-cohomology-]{sites-cohomology}
\externaldocument[dga-]{dga}
\externaldocument[dpa-]{dpa}
\externaldocument[sdga-]{sdga}
\externaldocument[hypercovering-]{hypercovering}
\externaldocument[schemes-]{schemes}
\externaldocument[constructions-]{constructions}
\externaldocument[properties-]{properties}
\externaldocument[morphisms-]{morphisms}
\externaldocument[coherent-]{coherent}
\externaldocument[divisors-]{divisors}
\externaldocument[limits-]{limits}
\externaldocument[varieties-]{varieties}
\externaldocument[topologies-]{topologies}
\externaldocument[descent-]{descent}
\externaldocument[perfect-]{perfect}
\externaldocument[more-morphisms-]{more-morphisms}
\externaldocument[flat-]{flat}
\externaldocument[groupoids-]{groupoids}
\externaldocument[more-groupoids-]{more-groupoids}
\externaldocument[etale-]{etale}
\externaldocument[chow-]{chow}
\externaldocument[intersection-]{intersection}
\externaldocument[pic-]{pic}
\externaldocument[weil-]{weil}
\externaldocument[adequate-]{adequate}
\externaldocument[dualizing-]{dualizing}
\externaldocument[duality-]{duality}
\externaldocument[discriminant-]{discriminant}
\externaldocument[derham-]{derham}
\externaldocument[local-cohomology-]{local-cohomology}
\externaldocument[algebraization-]{algebraization}
\externaldocument[curves-]{curves}
\externaldocument[resolve-]{resolve}
\externaldocument[models-]{models}
\externaldocument[functors-]{functors}
\externaldocument[equiv-]{equiv}
\externaldocument[pione-]{pione}
\externaldocument[etale-cohomology-]{etale-cohomology}
\externaldocument[proetale-]{proetale}
\externaldocument[relative-cycles-]{relative-cycles}
\externaldocument[more-etale-]{more-etale}
\externaldocument[trace-]{trace}
\externaldocument[crystalline-]{crystalline}
\externaldocument[spaces-]{spaces}
\externaldocument[spaces-properties-]{spaces-properties}
\externaldocument[spaces-morphisms-]{spaces-morphisms}
\externaldocument[decent-spaces-]{decent-spaces}
\externaldocument[spaces-cohomology-]{spaces-cohomology}
\externaldocument[spaces-limits-]{spaces-limits}
\externaldocument[spaces-divisors-]{spaces-divisors}
\externaldocument[spaces-over-fields-]{spaces-over-fields}
\externaldocument[spaces-topologies-]{spaces-topologies}
\externaldocument[spaces-descent-]{spaces-descent}
\externaldocument[spaces-perfect-]{spaces-perfect}
\externaldocument[spaces-more-morphisms-]{spaces-more-morphisms}
\externaldocument[spaces-flat-]{spaces-flat}
\externaldocument[spaces-groupoids-]{spaces-groupoids}
\externaldocument[spaces-more-groupoids-]{spaces-more-groupoids}
\externaldocument[bootstrap-]{bootstrap}
\externaldocument[spaces-pushouts-]{spaces-pushouts}
\externaldocument[spaces-chow-]{spaces-chow}
\externaldocument[groupoids-quotients-]{groupoids-quotients}
\externaldocument[spaces-more-cohomology-]{spaces-more-cohomology}
\externaldocument[spaces-simplicial-]{spaces-simplicial}
\externaldocument[spaces-duality-]{spaces-duality}
\externaldocument[formal-spaces-]{formal-spaces}
\externaldocument[restricted-]{restricted}
\externaldocument[spaces-resolve-]{spaces-resolve}
\externaldocument[formal-defos-]{formal-defos}
\externaldocument[defos-]{defos}
\externaldocument[cotangent-]{cotangent}
\externaldocument[examples-defos-]{examples-defos}
\externaldocument[algebraic-]{algebraic}
\externaldocument[examples-stacks-]{examples-stacks}
\externaldocument[stacks-sheaves-]{stacks-sheaves}
\externaldocument[criteria-]{criteria}
\externaldocument[artin-]{artin}
\externaldocument[quot-]{quot}
\externaldocument[stacks-properties-]{stacks-properties}
\externaldocument[stacks-morphisms-]{stacks-morphisms}
\externaldocument[stacks-limits-]{stacks-limits}
\externaldocument[stacks-cohomology-]{stacks-cohomology}
\externaldocument[stacks-perfect-]{stacks-perfect}
\externaldocument[stacks-introduction-]{stacks-introduction}
\externaldocument[stacks-more-morphisms-]{stacks-more-morphisms}
\externaldocument[stacks-geometry-]{stacks-geometry}
\externaldocument[moduli-]{moduli}
\externaldocument[moduli-curves-]{moduli-curves}
\externaldocument[examples-]{examples}
\externaldocument[exercises-]{exercises}
\externaldocument[guide-]{guide}
\externaldocument[desirables-]{desirables}
\externaldocument[coding-]{coding}
\externaldocument[obsolete-]{obsolete}
\externaldocument[fdl-]{fdl}
\externaldocument[index-]{index}

% Heap Project documents
\externaldocument[k3-]{k3}
\externaldocument[torelli-]{torelli}
\externaldocument[linear-algebra-]{linear-algebra}
\externaldocument[groups-]{groups}
\externaldocument[complex-analysis-]{complex-analysis}
\externaldocument[functional-analysis-]{functional-analysis}
\externaldocument[categories-]{categories}
\externaldocument[commutative-algebra-]{commutative-algebra}
\externaldocument[ringed-spaces-]{ringed-spaces}
\externaldocument[homological-algebra-]{homological-algebra}
\externaldocument[lie-algebras-]{lie-algebras}
\externaldocument[representation-theory-]{representation-theory}
\externaldocument[smooth-manifolds-]{smooth-manifolds}
\externaldocument[differential-topology-]{differential-topology}
\externaldocument[algebraic-topology-]{algebraic-topology}
\externaldocument[differential-geometry-]{differential-geometry}
\externaldocument[algebraic-geometry-]{algebraic-geometry}
\externaldocument[symplectic-geometry-]{symplectic-geometry}
\externaldocument[complex-geometry-]{complex-geometry}
\externaldocument[hodge-theory-]{hodge-theory}
\externaldocument[%
variations-of-hodge-structures-]{%
variations-of-hodge-structures}
\externaldocument[%
homological-mirror-symmetry-]{%
homological-mirror-symmetry}
\maybeexternaldocument{deformations-}{deformations}
\externaldocument[vertex-algebras-]{vertex-algebras}
\externaldocument[springer-]{springer}
\externaldocument[infty-categories-]{infty-categories}
\externaldocument[seminars-differential-geometry-]{%
seminars-differential-geometry}
\externaldocument[seminars-complex-geometry-]{%
seminars-complex-geometry}
\externaldocument[seminars-algebraic-geometry-]{%
seminars-algebraic-geometry}
\externaldocument[seminars-analysis-pde-]{%
seminars-analysis-pde}
\externaldocument[surfaces-]{surfaces}
\externaldocument[birational-maps-commutative-algebra-]{birational-maps-commutative-algebra}
\externaldocument[cremona-transformations-]{cremona-transformations}
\externaldocument[derived-categories-of-sheaves-]{derived-categories-of-sheaves}
\externaldocument[geometric-invariant-theory-]{geometric-invariant-theory}
\externaldocument[geometric-stability-conditions-and-group-actions-]{geometric-stability-conditions-and-group-actions}
\externaldocument[moduli-spaces-of-sheaves-]{moduli-spaces-of-sheaves}
\externaldocument[bridgeland-stability-general-theory-]{bridgeland-stability-general-theory}
\externaldocument[hilbert-schemes-]{hilbert-schemes}
\externaldocument[noncommutative-abelian-surfaces-and-kummer-type-hyperkahler-varieties-]{noncommutative-abelian-surfaces-and-kummer-type-hyperkahler-varieties}
\externaldocument[mumford-tate-groups-in-hodge-theory-]{mumford-tate-groups-in-hodge-theory}
\externaldocument[hodge-birational-atoms-2026-]{hodge-birational-atoms-2026}
\externaldocument[xxii-egd-2026-]{%
xxii-egd-2026}
\externaldocument[pde-]{pde}
\externaldocument[probability-]{probability}

\externaldocument[linguistics-]{linguistics}

\externaldocument[%
16th-alga-meeting-2026-]{%
16th-alga-meeting-2026}

\externaldocument[%
iv-jornadas-lefschetz-2026-]{%
iv-jornadas-lefschetz-2026}

\externaldocument[reading-the-masters-]{%
reading-the-masters}

\externaldocument[real-analysis-]{%
real-analysis}

\externaldocument[optimal-transport-]{%
optimal-transport}

\externaldocument[history-]{history}

\externaldocument[seiberg-witten-]{seiberg-witten}

\externaldocument[history-philosophy-biology-]{history-philosophy-biology}

\externaldocument[classical-mechanics-]{classical-mechanics}
\externaldocument[electromagnetism-]{electromagnetism}
\externaldocument[waves-and-optics-]{waves-and-optics}
\externaldocument[wif-]{waves-in-fluids-iv}
\externaldocument[thermal-physics-]{thermal-physics}
\externaldocument[quantum-mechanics-]{quantum-mechanics}
\externaldocument[special-relativity-]{special-relativity}
\externaldocument[physics-]{physics}

\externaldocument[chemistry-]{chemistry}

\externaldocument[molecular-cell-biology-]{molecular-cell-biology}

\externaldocument[microbiology-]{microbiology}

\externaldocument[protistology-]{protistology}

\externaldocument[mycology-]{mycology}

\externaldocument[botany-]{botany}

\externaldocument[zoology-]{zoology}

\externaldocument[animal-histology-]{animal-histology}

\externaldocument[animal-physiology-]{animal-physiology}

\externaldocument[systematics-biogeography-]{systematics-biogeography}

\externaldocument[ecology-natural-resources-]{ecology-natural-resources}

\externaldocument[earth-sciences-]{earth-sciences}

\externaldocument[biostatistics-]{biostatistics}

\externaldocument[mathematical-systems-biology-]{mathematical-systems-biology}

\externaldocument[cimpa-el-salvador-2024-]{cimpa-el-salvador-2024}

\externaldocument[anthropology-]{anthropology}

\externaldocument[yang-mills-theory-]{yang-mills-theory}

\externaldocument[seminars-algebra-]{seminars-algebra}

\externaldocument[seminars-physics-]{seminars-physics}

\externaldocument[seminars-computing-]{seminars-computing}

\externaldocument[oktobermat-2026-]{oktobermat-2026}

\externaldocument[saet-2026-]{saet-2026}

\externaldocument[seminars-symplectic-geometry-]{seminars-symplectic-geometry}

\externaldocument[seminars-dynamics-]{seminars-dynamics}

\externaldocument[oktobermat-2024-]{oktobermat-2024}

\externaldocument[quantum-groups-]{quantum-groups}

\externaldocument[legal-2027-]{legal-2027}

\externaldocument[derived-categories-]{derived-categories}

% Theorem environments.
%
\theoremstyle{plain}
\newtheorem{theorem}[subsection]{Theorem}
\newtheorem{proposition}[subsection]{Proposition}
\newtheorem{lemma}[subsection]{Lemma}
\newtheorem{corollary}[subsection]{Corollary}

\theoremstyle{definition}
\newtheorem{definition}[subsection]{Definition}
\newtheorem{example}[subsection]{Example}
\newtheorem{exercise}[subsection]{Exercise}
\newtheorem{situation}[subsection]{Situation}

\theoremstyle{remark}
\newtheorem{remark}[subsection]{Remark}
\newtheorem{remarks}[subsection]{Remarks}

\numberwithin{equation}{subsection}

% Macros
%
\def\lim{\mathop{\mathrm{lim}}\nolimits}
\def\colim{\mathop{\mathrm{colim}}\nolimits}
\def\Spec{\mathop{\mathrm{Spec}}}
\def\Hom{\mathop{\mathrm{Hom}}\nolimits}
\def\Ext{\mathop{\mathrm{Ext}}\nolimits}
\def\SheafHom{\mathop{\mathcal{H}\!\mathit{om}}\nolimits}
\def\SheafExt{\mathop{\mathcal{E}\!\mathit{xt}}\nolimits}
\def\Sch{\mathit{Sch}}
\def\Mor{\mathop{\mathrm{Mor}}\nolimits}
\def\Ob{\mathop{\mathrm{Ob}}\nolimits}
\def\Sh{\mathop{\mathit{Sh}}\nolimits}
\def\NL{\mathop{N\!L}\nolimits}
\def\CH{\mathop{\mathrm{CH}}\nolimits}
\def\proetale{{pro\text{-}\acute{e}tale}}
\def\etale{{\acute{e}tale}}
\def\QCoh{\mathit{QCoh}}
\def\Ker{\mathop{\mathrm{Ker}}}
\def\Im{\mathop{\mathrm{Im}}}
\def\Coker{\mathop{\mathrm{Coker}}}
\def\Coim{\mathop{\mathrm{Coim}}}

% Boxtimes
%
\DeclareMathSymbol{\boxtimes}{\mathbin}{AMSa}{"02}

%
% Macros for moduli stacks/spaces
%
\def\QCohstack{\mathcal{QC}\!\mathit{oh}}
\def\Cohstack{\mathcal{C}\!\mathit{oh}}
\def\Spacesstack{\mathcal{S}\!\mathit{paces}}
\def\Quotfunctor{\mathrm{Quot}}
\def\Hilbfunctor{\mathrm{Hilb}}
\def\Curvesstack{\mathcal{C}\!\mathit{urves}}
\def\Polarizedstack{\mathcal{P}\!\mathit{olarized}}
\def\Complexesstack{\mathcal{C}\!\mathit{omplexes}}
% \Pic is the operator that assigns to X its picard group, usage \Pic(X)
% \Picardstack_{X/B} denotes the Picard stack of X over B
% \Picardfunctor_{X/B} denotes the Picard functor of X over B
\def\Pic{\mathop{\mathrm{Pic}}\nolimits}
\def\Picardstack{\mathcal{P}\!\mathit{ic}}
\def\Picardfunctor{\mathrm{Pic}}
\def\Deformationcategory{\mathcal{D}\!\mathit{ef}}


\begin{document}

\title{Seminars in Dynamics}
\maketitle

\phantomsection
\label{section-phantom}

\tableofcontents



\section{Medidas estacionárias de um Sistema
Dinâmico Aleatório no círculo}
\label{section-medidas-estacionarias-circulo}
%% speaker: Graccyela Rosybell Salcedo
%% speaker: Pirela, ICMC
%% place: Ergodic Theory Seminar, IMPA
%% date: September 26, 2024

\noindent
Graccyela Rosybell Salcedo Pirela, ICMC.
Ergodic Theory Seminar, IMPA.
26 de setembro de 2024.

\medskip\noindent{\bf Exemplos.}
Rotações  com pendente racional ou
irracional. Sistema com pontos atratores e
repulsores.

\medskip\noindent{\bf Conf.}
$\mathcal{F}\subset
\operatorname{Hom}(\mathbb{S}^1)$,
$\nu$ probabilidade,
$\operatorname{Hom}(\mathbb{S}^1)$,
$\operatorname{supp}(\nu)=\mathcal{F}$.
Espaço de probabilidade:
$$
(\Omega,\mathbb{P})=
(\mathcal{F}^{\mathbb{N}},\nu^{\mathbb{N}})
$$
Caminho aleatório:
$$
\Omega \ni \omega\longmapsto
(f^n_\omega)_{n\geq 0}
$$
onde $\omega=(f_n)_{ n \in \mathbb{N}}$,
$f_{\omega}^n:=f_n\circ \ldots \circ f_1$,
$f^0_\omega =
\operatorname{id}_{\mathbb{S}^1}$.

\medskip\noindent{\bf Ponto de vista
topológico.}
$$
S=\{f^n_{\omega}:\omega \in\Omega,n \geq 0\}
$$
$A\subset \mathbb{S}$ é
\textit{invariante} se
$f(A)\subset A$, $\forall f\in S$.
Um conjunto fechado
$A\subset \mathbb{S}^1$ é chamado
\textit{minimal} invariante se
$\forall B\subset A$ tal que $B$ é
invariante então $B=A$ ou  $B=\varnothing $.
$$
O(x)=\{f(x):f\in S\}
$$

\medskip\noindent{\bf Ponto de vista
probabilístico.}
$(f^n_\omega(x))$.
$\eta\in\operatorname{Prob}(\mathbb{S}^1)$,
dizemos que é  \textit{$\nu$-estacionária} se
$$
\eta=\int f_*\eta d\nu(f).
$$
$$
T:(\omega,x)\longmapsto
(\sigma\omega,f'_\omega(x))
$$
$\eta$ é $\nu$-estacionária  $\iff$
$\mathbb{P}\otimes \eta$ e $T$-invariante.
Dizemos que $\eta$ é
\textit{$\nu$-ergódica e estacionaria} se
$\mathbb{P}\otimes \eta$ é $T$-ergódica.

\medskip\noindent{\bf Resultados previos.}

\begin{theorem}[Resultado previo; Furst, 63]
\label{theorem-sp-furst}
$\forall x:\mathbb{P}$ $\omega$
$$
d\left( \frac{1}{N}\sum_{n=0}^{N-1}
\delta_{p^n_\omega x} \right)
$$
$\subset$ $\nu$-estacionária
\end{theorem}

\begin{theorem}[Resultado previo;
Malicet '17]
\label{theorem-sp-malicet}
Tricotomía:
\begin{enumerate}
\item $\not \exists
m\in\operatorname{Prob}(\mathbb{S}^1)$,
$f_*m=m\;\forall f\in\mathcal{F}\implies $
contração local $\implies $ $\exists
d\in\mathbb{N}$,
$\operatorname{Prob}_{\nu}^
{\operatorname{erg}}
(\mathbb{S}^1)=\{\eta_1,\ldots,\eta_d\}$
$$
k:=\operatorname{supp}(\eta)
$$
$$
\implies U_i(x)=\mathbb{P}\left(\omega
\frac{1}{N}\sum_{n=0}^{N-1}
\delta_{p_\omega^\eta}\right)
\overset{\omega}{\longrightarrow}\eta_1
$$

\item O semigrupo $\delta$ é topológico
conjugado como um semigrupo das isometrías
atuando minimalmente
($\mathbb{S}^1$ é mínima).

$\exists !$ medida $\nu$-esacionária

\item $\exists x$
$O(x)=\{f(x),\ldots,f\in\mathbb{S}\}$
infinita.
\end{enumerate}
\end{theorem}

\begin{example}[Exemplos]
\label{example-sp-exemplos}
\begin{enumerate}
\item $f_1=R_{p/q}$.
\item Círculo con atrator no polo norte e
repulsor no polo sul.
\item[Desenho] Parece que são círculos com
dois pontos mais o fluxo sai de cada ponto a
um lado dele e entra do outro lado.
\end{enumerate}
\end{example}

\medskip\noindent{\bf $|O(x)|= \infty\;
\forall x$.}
$$
m^* _x=\lim_{N\to \infty}\frac{1}{N}
\sum_{n=0}^{N-1}\int
\delta_{\rho^\eta_\omega x}
d\mathbb{P}(\omega)
=\sum_{i=1}^d U_i(x)\eta_i
$$
$x\longrightarrow m_x$ contínua.

\begin{itemize}
\item[$d=1$] $x\longrightarrow \eta_1$.
\item[$d=2$] $x\longrightarrow m_x$
\item[$d>2$]
[Desnhos de triangulo e quadrado]
\end{itemize}

\begin{theorem}[Teorema 1]
\label{theorem-sp-1}
$\forall x\; |O(x)|= \infty$
$\forall i=1,\ldots,d$ $ \exists
\mathcal{A}_{i}$ família finita de
intervalos fechados tal que
\begin{enumerate}
\item $K_1\subset \{U=1\} =
\bigcup_{I\in A_i}I $

\item $i\neq j$,
$I\cap J=\varnothing,\;\forall J\in A_i,
J\in A_j$.

\item $\forall a,b,$ $a\in J\in A_i,$
$b \in J\in A_j$ tal que
$(a,b)\cap \bigcup_{i=1} ^d
\bigcup_{J\in A_i} I\neq \varnothing$.
\begin{enumerate}
\item[(i)] $U_k=0$ em $[a,b]$
$\forall k\neq 1,0$.
\item[(ii)] $U_i,U_j$ são monótonos sobre
$(a,b)$.
\end{enumerate}
\end{enumerate}
\end{theorem}

\begin{corollary}
\label{corollary-sp-orbit}
$\forall x$ $|O(x)| =\infty$,
$\{\eta_x:x\in\mathbb{S}^1\} \subset
\{t\eta_1+(1-t)\eta_j:N\in\{1,\ldots,d\},
t\in [0,1]\}$.
\end{corollary}

\begin{proposition}[Proposição 2]
\label{proposition-sp-2}
$I\subset \mathbb{S}^1$ intervalo,
$i\in\{1,\ldots,d\}$.
Se $U_1|_{I}\equiv C_f$
\end{proposition}

\begin{proof}[Do teorema 1]

\end{proof}

\medskip\noindent
{\bf Editorial note.}
The original notes are unfinished and
transcribed during the talk: the statement of
Proposition~\ref{proposition-sp-2}
stops in the middle of a sentence, the proof
of Theorem~\ref{theorem-sp-1} is empty, and
several formulas contain evident slips
(for instance the symbols $p$, $\rho$
and $f'$ in place of $f$, the index $N$ in
the corollary, and $\mathbb{S}$ for
$\mathbb{S}^1$). They are kept as in the
original. Changes made: the original
macros for theorem-like environments became
the standard environments, the enumerate
labels became explicit item labels, the
text inside two displays was moved to the
surrounding lines, and the month
``Septembro'' was written ``setembro''.

\section{Ação de homeos sobre o grafo fino de
curvas: uma classificação}
\label{section-acao-homeos-grafo-fino}
%% speaker: Pierre-Antoine Guihéneuf,
%% speaker: Sorbonne Université
%% place: Ergodic Theory Seminar, IMPA
%% date: October 24, 2024

\noindent
Pierre-Antoine Guihéneuf,
Sorbonne Université.
Ergodic Theory Seminar, IMPA.
24 de Outubro de 2024.

\medskip\noindent{\bf Grafo fino de curvas.}
Fixe $S$ uma superfície fechada de género
$g \geq 1$. O \textit{grafo fino de curvas}
$C^+(S)$ (tem um grafo que não é fino que é
exatamente a mesma coisa só que
\textit{up to homotopy}) é:
\begin{itemize}
\item Vértices: curvas simples fechadas não
homotopicamente triviais em $S$.
\item Arestas: $\alpha-\beta$ sse
$\#(\alpha\cap \beta)\leq 1$.
\end{itemize}
Tem uma distância natural em $C^+(S)$.

\begin{remark}
\label{remark-gh-isometries}
$\operatorname{Homeo}(S)\curvearrowright
C^+(S)$ dada por $h\alpha=h(\alpha)$ é uma
ação por isometrias.
\end{remark}

\begin{remark}
\label{remark-gh-not-compact}
$C^+(S)$ não é localmente compacto.
\end{remark}

\begin{theorem}[BHW21 $\ddot\smile$]
\label{theorem-gh-bhw21}
$C^+(S)$ é conexo, de diametro infinito e
Gromov-hiperbólico.
\end{theorem}

\begin{definition}
\label{definition-gh-hyperbolic}
Um grafo é \textit{Gromov-hiperbólico} se
$\exists \delta>0$ tal que todos os
triangulos geodésicos são $\delta$-finos.
\end{definition}

\begin{definition}
\label{definition-gh-thin-triangle}
\textit{Triangulo $\delta$-fino}: que pode
pegar uma vizinhança de raio $\delta $ de
qualquer par de lados, e o lado restante
fica conteúdo na união daquelas
vizinhanças.
\end{definition}

\medskip\noindent{\bf Um pouco sobre espaços
Gromov-hiperbólicos.}

\begin{example}
\label{example-gh-hyperbolic}
\begin{itemize}
\item Arvores
\item $\mathbb{H}^2$
\end{itemize}
\end{example}

\medskip\noindent{\bf Classificação das
isometrias.}
\begin{itemize}
\item \textit{Eliptica} se existe uma órbita
de diametro finito, i.e. $\exists x \in X$
tal que
$\operatorname{diam}\{f^n(x)\}<+\infty$.
\item \textit{Loxodrómica} se
$\exists x \in X$ tal que
$$
\lim_{n \to \infty}
\frac{d(f^n(x),x)}{n}
$$
\item \textit{Parabólica} o resto.
\end{itemize}

\begin{definition}
\label{definition-gh-quasi-isometry}
Uma \textit{quase-isometria} de $X$ é
$g:X\to X$ tal que $\exists \lambda,c$
$$
\lambda^{-1}d(x,y)-c\leq d(g(x),g(y))
\leq \lambda d(x,y)+c
$$
\end{definition}

\medskip\noindent{\bf O bordo de $X$.}

\begin{definition}
\label{definition-gh-boundary}
O \textit{bordo de $X$}, denotado
$\partial X$, é o conjunto dos mergulhos
$Q$ de $\mathbb{R}_+$ em $X$ módulo
distância de Hausdorff finita, isto é, o
conjunto de direções no infinito.
\end{definition}

\begin{proposition}
\label{proposition-gh-complete}
$X\cup  \partial X$ é completo.
\end{proposition}

\medskip\noindent{\bf Pergunta.}
Forma de $\partial C^+(S)$?

\begin{remark}
\label{remark-gh-extends}
Toda isometria de $X$ se extende em um
homeo de $X\cup \partial X$.
\end{remark}

\begin{theorem}[Gromov]
\label{theorem-gh-gromov}
\begin{itemize}
\item Se $f$ é parabólica então $f$ tem um
único ponto fixo em $X\cup \partial X$, qué
esta em $\partial X$.

\item Se $f$ é loxodrómico, então $f$ tem
dois pontos fixos em $X \cup \partial X$ que
estão no bordo e a dinámica é norte-sul.
$f\Big|_{\partial X\setminus \{a^-\}}$ é
uma contração.
\end{itemize}
\end{theorem}

\begin{remark}[Dani]
\label{remark-gh-dani}
Parece que tem uma dinâmica dada simplesmente
pela iteração da $f$. Supongo que essa $f$ é
um isomorfismo.
\end{remark}

$f^{-1}|_{\partial X\setminus \{a^+\}}$ é
uma contração. Isso permete jogar ping-pong.

\begin{exercise}
\label{exercise-gh-free}
$\exists n$ tal que
$\left\langle{} f^n,g^n \right\rangle{}$
é livre.
\end{exercise}

\begin{theorem}[Hensel, le Raux (private
communication)]
\label{theorem-gh-hensel-le-raux}
Se $f,g\in\operatorname{Homeo}(\Pi^2)$,
$\rho(f),\rho(g)$ com interior não vazio,
$\rho(f),\rho(g)$ não são os mesmos ao menos
de homotecia/traslação. Então existe $n$ tal
que
$\left\langle{} f^n,g^n \right\rangle{}$
é livre e $\forall h\in
\left\langle{} f^n,g^n \right\rangle{}$,
$\operatorname{int}(\rho(h))\neq
\varnothing$.
\end{theorem}

\medskip\noindent{\bf Classificação.}

\begin{theorem}
\label{theorem-gh-classification}
$f\in\operatorname{Homeo}_0(S)$. Então aspse
\begin{enumerate}
\item[(i)] $f$ age sobre $C^+(S)$ de maneira
loxodromica.
\item[(ii)] $\exists P\subset S$ finito,
$f$-invariante tal que $f|_{S\setminus P}$ é
pseudo-Anosov.
\item[(iii)]
$\operatorname{int}\rho_{\operatorname{erg}}
(f)\neq \varnothing$.
\end{enumerate}
\end{theorem}

\begin{theorem}
\label{theorem-gh-torus}
$f\in\operatorname{Homeo}_0(\Pi^2)$.
Então $f$ age de maneira elíptica see $f$
tem desvios limitados em uma direção
racional $v\in\mathbb{Z}^2\setminus\{0\}$,
$\exists C,p\in\mathbb{R}^{2}$ tais que
$\forall n,x$,
$$
|\left\langle{} \tilde{f}^n(x)-n\rho,v
\right\rangle{} |\leq C
$$
\end{theorem}

\medskip\noindent
{\bf Editorial note.}
The original text is kept as written,
including its apparent slips (``aspse'' in
Theorem~\ref{theorem-gh-classification},
the limit in the definition of loxodromic
isometries stated without its condition,
$\Pi^2$ and the unexplained $\rho$ and $p$).
Changes made: the definition of the
boundary $\partial X$ was an aligned display
containing text and is now a sentence with
the same words; broken operator names
(``di a m'', ``i n t'', ``Ho m eo'') were
repaired; inner products were written in the
form required by the notebook style; the
term defined in
Definition~\ref{definition-gh-thin-triangle}
was italicized; the boxes ``Classificação
das isometrias'' and ``Pergunta'' and the
empty box ``O bordo de $X$'' became bold
run-in headings.

\section{Existence and non-existence of
physical measures for doubly intermittent
maps}
\label{section-physical-measures-%
intermittent}
%% speaker: Stefano Luzzatto, ICTP
%% place: Ergodic Theory Seminar, IMPA
%% date: November 28, 2024

\noindent
Stefano Luzzatto, ICTP.
Ergodic Theory Seminar, IMPA.
November 28, 2024.

\medskip\noindent
{\bf Abstract.}
We introduce a large class of one-dimensional
map in the topological conjugacy class of
2x mod1 but exhibiting a variety of ergodic
behaviours, such as the existence of
invariant probability measures equivalent to
Lebesgue, Dirac-delta physical measures and,
perhaps most interestingly, non-existence of
physical measure. This class generalises the
standard well known and well studied
Pomeau-Manneville and
Liverani-Saussol-Vaienti maps.

\medskip\noindent{\bf Physical measures.}
$X$ compact measure space, $m$ a reference
measure, $f:X \to X$, $x \in X$.
$$
e_n(x):=\frac{1}{n}\sum_{i=0}^{n-1}
\delta_{f^i(x)}
$$
é a \textit{sequência canônica de empirical
measures}.

Se $e_n(x)$ converge a $\mu$. Então $\mu$
descreve a \textit{estadística}  de $x$. O
\textit{basin} de $\mu$ $\mathcal{B}_\mu$ é o
conjunto de pontos $x \in X$ tal que
$e_n(x)\to \mu$. $\mu$ é uma
\textit{physical measure} se
$m(\mathcal{B}_{\mu})>0$.

\medskip\noindent{\bf Pergunta.}
O que acontece se $e_n(x) $ não converge?

Devem existir duas medidas $\mu,\nu$ e
$n_i,n_j\to\infty$ tais que
$e_{n_i}(x)\to \mu$ e $e_{n_j}\to \nu$.
Dizemos que $x$ tem um comportamento
\textit{não estadístico}. Se $m$  é tal que
quase todo ponto é não estadístico, $m$ é
\textit{não estadística}.

\medskip\noindent{\bf Exemplos.}
Bota muitos zeros, depois muito mais uns,
depois muuuito mais zeros…

\medskip\noindent
{\bf Editorial note.}
The original notes stop after the first
example and contain nothing about the
intermittent maps of the title. Repairs: the
summation index in the definition of $e_n(x)$
was written $n=0$ (now $i=0$), and the
names ``PomeauManneville'' and
``Liverani-Saussol0-Vaienti'' in the
abstract, garbled in the original, were
written ``Pomeau-Manneville'' and
``Liverani-Saussol-Vaienti''. The
speaker's name, spelled ``Luzzato'' in the
original header, is written Luzzatto.

\section{Substitution dynamics on infinite
alphabets}
\label{section-substitution-dynamics}
%% speaker: Ali Messaoudi, UNESP
%% place: Ergodic Theory Seminar, IMPA
%% date: November 28, 2024

\noindent
Ali Messaoudi, UNESP.
Ergodic Theory Seminar, IMPA.
November 28, 2024.

\medskip\noindent
{\bf Abstract.}
A substitution is a map from an alphabet A to
the set of finite words in A. To any
substitution we can naturally associate a
symbolic dynamical system that is well
studied in the literature when the alphabet
is a finite set and connected to several
areas such as ergodic theory and number
theory among others.
In this work, we study ergodic and geometric
properties of dynamical systems associated to
substitutions in infinite alphabets. This
study involves finite and infinite invariant
measures, countably infinite matrices and
Rauzy Fractals.

\medskip\noindent{\bf Introdução para quem
não é da área.}

\begin{definition}
\label{definition-ms-dynamical-system}
Um \textit{sistema dinâmico} é um par
$(X,f)$ com  $X \neq  \varnothing$ and
$f:X\to X$.
\end{definition}

Em sistemas dinâmicos estamos interessados em
estudar a órbita
$$
O(x)=\{f^n(x), n \in \mathbb{N},
f^n:=f \circ f \circ\ldots \circ f\}
$$
para $x \in X$.

Quando tem uma medida invariante, i.e.
$\mu(f^{-1}(A))=\mu(A),\; \forall A \in
\mathcal{B}$,  quase todo ponto volta perto
dele mesmo.

\medskip\noindent{\bf Introdução de verdade.}
Considere um alfabeto $A$ que é só um
conjunto finito. $A^*$ são as palavras
finitas. $A^\mathbb{N}$ são as palavras
infinitas. $A^\mathbb{Z}$ são as palavras
infinitas pra diante e pra atrais.

Uma \textit{substitução} é um mapa
$$
\begin{aligned}
\sigma: A &\longrightarrow A^* \\
a &\longmapsto \sigma(a)
\end{aligned}
$$

\begin{lemma}
\label{lemma-ms-periodic}
Existe $u \in A^\mathbb{N}$ e
$k \in \mathbb{N}^*$ tal que
$\sigma^k(u)=u$ é um ponto periódico de
$\sigma$.
\end{lemma}

\begin{remark}
\label{remark-ms-fibonacci}
A substituição de Fibonacci $A=\{0,1\}$
$\sigma(0)=01, \sigma(1)=0$ tem um ponto
fixo, o limite de $\sigma(0)$, ie.
$$
u:=\lim_{n \to \infty} \sigma^n(0),\qquad
\sigma(u)=u.
$$
Pode codificar a informação da substituição
$\sigma$ numa matriz: cada renglón é o numero
de vezes que aparece cada letra do alfabeto
na imagem de $\sigma$.
\end{remark}

\medskip\noindent{\bf Classical results.}

\medskip\noindent
{\bf Editorial note.}
The original notes end with the heading
``Classical results'' and no text below it;
nothing was added. Repairs: a missing
parenthesis in $\mu(f^{-1}(A))=\mu(A)$
and the stray ``,;'' before $\forall A$
were fixed, and
the aligned display of $\sigma$ no longer
contains text. The speaker's name, spelled
``Messãoudi'' in the original header, is
written Messaoudi.

\section{Oseledets theorem}
\label{section-oseledets-theorem}
%% speaker: Dimitrii Korshunov, IMPA
%% place: Geometric Structures Seminar, IMPA
%% date: February 27, 2025

\noindent
Dimitrii Korshunov, IMPA.
Geometric Structures Seminar.
February 27, 2025.

\medskip\noindent
{\bf Abstract.}
Risking to sell coal to Newcastle, we will
prove the multiplicative ergodic theorem:
given a metrized bundle E over an ergodic
dynamical system (M,T) and a lifting of T to
an action on E (i.e. a cocycle), there is a
filtration of E by subbundles $E_i$ such that
lengths of iterates of vectors from
$E_i\setminus E_{i+1}$ have the same
asymptotic behavior and grow like
$\operatorname{exp}{\lambda_i N}$,
$\lambda_i$ being the famous Lyapunov
exponents. Following Kaimanovich, we will
emphasise the geometric meaning of this
theorem, centered around the notion of
regular sequence in symmetric spaces. If time
permits, we will discuss how this story can
be thought of as the ``law of large numbers
for semi-simple Lie groups''.

\medskip\noindent{\bf Introduction.}
Setting: an ergodic dynamical system
$(\Omega,\sigma,\mu,T)$ where $\Omega$ is a
separable, second countable metric space,
$\sigma$ is a Borel sigma-algebra and $\mu$
measure, and $T$ measure-preserving.

\begin{definition}
\label{definition-ko-ergodic}
$T$ is \textit{ergodic} if for $U$
measureable subset of  $\Omega$ with
$T(U) \subset U$ we have that $\mu(U) =0$ or
$\mu(U)=1$.
\end{definition}

\begin{theorem}[Birkhoff ergodic theorem]
\label{theorem-ko-birkhoff}
Idea: that time average is spatial average.

For ergodic system, if $f=f^+ + f^-$ and
$\int_\Omega f^+ < \infty$,
$$
\lim_{N \to \infty}
\sum_{i=1}^{N-1}f(T^i \omega)
=\int_\Omega fd\mu
$$
for  $\mu$-almost any $\omega$.
\end{theorem}

\begin{definition}
\label{definition-ko-cocycle}
\textit{Cocycle} is morally the following
pair: a smooth function on a manifold to
itself and its differential, so cocycle is
$(f,df)$. More abstractly it will be some map
of the fibers of a vector bundle over
$\Omega$. So a map of the total spaces of the
bundle that restricted to fibers is linear.
\end{definition}

\medskip\noindent{\bf Oseledets theorem.}

\begin{theorem}[Oseledets most probably]
\label{theorem-ko-oseledets}
Given a cocycle over ergodic dynamical system
and suppose that $V_\omega$ is equipped with
$\|\cdot \|$ such that
$$
\int_\Omega \operatorname{log}^+
\|T\|_{\operatorname{op}}d\mu(\omega)
< \infty
$$
where $\operatorname{log}_+ f=0$ if
$\operatorname{log}f<0$ and
$\operatorname{log}_+ f=\operatorname{log}f$
otherwise. Then there exists a set oreal
numbers
$\lambda_1>\lambda_2>\ldots>\lambda_{n-1}$
and an invariant flag (=nested vector spaces)
$$
V=C^{\leq  \lambda_1}\supset
V^{\leq  \lambda_2}\supset \ldots\supset
V^{\leq  \lambda_{n-1}}
$$
whwere $\lambda_{n-1}$ could be $-\infty$.
such that $TV^{\leq \lambda_i}\subset
V^{\leq  \lambda_i}$ and for any
$v \in V^{\leq  \lambda_i}\setminus
V^{\leq \lambda_{i+1}}$ then
(idea: $Tv$ grows exponentially with exponent
$\lambda_i$)
$$
\lim_{N} \frac{1}{N}\operatorname{log}
\| T^Nv\|=\lambda_i
$$
\end{theorem}

The $\lambda_i$ are called
\textit{Lyapunov exponents} and they have a
lot information about the dynamical system.

\medskip\noindent{\bf Birkhoff to
1-dimensional Oseledets.}
Let's apply Birkhoff on
$$
f(\omega)=\operatorname{log}
\frac{\|T_\omega v\|}{\| v\|}
$$
(Notice that we can apply Birkhoff because
the fact that $\int_\Omega f^+ d\mu< \infty$
is precisely the  condition on Oseledets of
$\operatorname{log}^+$)

Anyway we compute that the sum is the
logarithm of a telescopic product, and by
Birkhoff
$$
\frac{1}{N}\Big(f(\omega)+\ldots+
f(T^{N-1}\omega)\Big)
=\frac{1}{N}\operatorname{log}
\frac{\|T^Nv\|}{\|v\|}
\to\int_\Omega f d \mu(\omega)=\lambda_1
$$

\medskip\noindent{\bf 1-dimensional Oseledets
to Birkhoff.}
Oops! But you take the exponent of the
function and apply Oseledets.

\medskip\noindent
{\bf Structure of the proof.}

\begin{lemma}[Slogan: reducibility or growth
dicotomy]
\label{lemma-ko-dichotomy}
One of the following holds (actually both can
hold)
\begin{enumerate}
\item There is just one Lyapunov exponent
(any vector grows at the same rate), i.e.
$$
\forall v\qquad \lim_{N \to \infty}
\operatorname{log}\|T^N v\|=\lambda_1
$$
\item there is a proper invariant subbundle
$E \subset V$ with
$0<\dim E< \dim V$.
\end{enumerate}
\end{lemma}

\begin{lemma}[Slogan: unusual growth implies
splitting]
\label{lemma-ko-splitting}
$$
0\to E\to V\to F\to 0
$$
[Figure in the original notes: a commutative
diagram with this short exact sequence and a
dashed arrow $\sigma$ bending back from
$F$ to $V$.]

Suppose there exists $\lambda_E$ with
$\operatorname{log} \|T^N e\|\to \lambda_F$
such that
$$
\forall e \in E\setminus\{0\},\qquad
\frac{1}{N}\operatorname{log}\|T^Ne\|
\to \lambda_F
$$
and
$$
\forall f \in E\setminus\{0\},\qquad
\lim_{N\to \infty}
\frac{1}{N}\operatorname{log}\|T^N f\|
\to \lambda_F
$$
and limits are uniform $\|e\|=1,\|f\|=1$.

If $\lambda_E > \lambda_F$ then the sequence
splits, that is, there exists
$\sigma:F \to E$ that is  $T$-invariant and
$$
V=E \oplus \sigma(F)\qquad \text{and}
\qquad p\cdot \sigma=\operatorname{Id}.
$$
\end{lemma}

\begin{definition}
\label{definition-ko-tempered}
\textit{Tempered}
$$
\lim_{N\to \infty}
\frac{1}{N}\operatorname{log}
\frac{\|\sigma(T^N v)\|}{\|T^Nv\|}=0.
$$
\end{definition}

We can prove the theorem modulo these two
lemmas.

\begin{proof}
Sketched
\end{proof}

\medskip\noindent
{\bf The geometric meaning of this
theorem.}

\begin{definition}
\label{definition-ko-regular}
A sequence is called \textit{regular} if
there exists a geodesic $\gamma$ and
$\theta>0$ such that
$$
\lim \frac{1}{n}d(p_i,
\gamma(\theta\cdot i))=c
$$
\end{definition}

\begin{remark}
\label{remark-ko-symmetric}
This is supposed to be studied when the
fibers of the bundle are symmetric spaces.
The symmetric spaces are the spaces of
metrics. I think this is the twistor space.
\end{remark}

\medskip\noindent
{\bf Editorial note.}
The section title and the text spell the
speaker's theorem ``Oseldets'' and
``Lyapunv''; both were written Oseledets and
Lyapunov. The proof is only announced
(``Sketched'') in the original. Formulas
repaired: the missing parenthesis in the
setting, the garbled braces in
$V^{\leq \lambda_i}\setminus
V^{\leq \lambda_{i+1}}$, a text case
distinction for $\operatorname{log}_+$ (now
a sentence), and the telescoping computation,
which in the original was an aligned display
with text and had the factor $1/N$ missing
on its first line. The apparent slips
below were left as written: $C^{\leq
\lambda_1}$ for $V$ in the flag, the Birkhoff
sum without $1/N$, the limit in
Lemma~\ref{lemma-ko-dichotomy} without $1/N$,
and in Lemma~\ref{lemma-ko-splitting} the
exponents $\lambda_E$ and $\lambda_F$ and the
vectors ``$f \in E$''. The commutative
diagram of Lemma~\ref{lemma-ko-splitting}
was a tikz-cd picture; its row is
written as a formula and the dashed arrow
is described in brackets. The original
header of this talk gives the series as
Geometric Structures Seminar.


\bibliography{my}
\bibliographystyle{amsalpha}

\end{document}
